\setlength{\absparsep}{18pt} 

\begin{resumo}[Resumo]
 \begin{otherlanguage*}{brazil}
   
   \vspace{\onelineskip}
   \noindent 
   
Estudantes da Educação Básica convivem com aplicações de inteligência artificial, mas esse contato não assegura a compreensão dos dados, regras e escolhas envolvidos em seu funcionamento. Este projeto de qualificação tem como objetivo desenvolver e analisar uma sequência didática, apoiada por uma plataforma interativa, para o ensino introdutório de \textit{machine learning} a estudantes do Ensino Médio. O tema é tratado como conteúdo de Computação, com atenção ao modo como modelos computacionais utilizam dados para produzir previsões, classificações e agrupamentos. A proposta pedagógica fundamenta-se nas normas curriculares da Resolução CNE/CEB nº 1/2022, no pensamento computacional e no aporte de múltiplas representações. As atividades seguem o ciclo prever, experimentar, observar e explicar. O produto educacional articula uma sequência didática de oito encontros a uma plataforma web interativa executada localmente com dados sintéticos tabulares. A sequência contempla regressão linear, regressão logística, \textit{k}-vizinhos mais próximos (\textit{k}-NN), árvore de decisão e \textit{k}-means. A pesquisa é de natureza aplicada, com abordagem predominantemente qualitativa e apoio de dados quantitativos descritivos, e será desenvolvida como intervenção pedagógica em uma turma de Ensino Médio sob docência do pesquisador. A produção de dados reunirá diagnóstico individual, registros individuais breves em três encontros, instrumento final com questões comparáveis às iniciais, desafio de transferência e registros docentes. A análise examinará mudanças, permanências e dificuldades em quatro dimensões: reconhecimento de dados e finalidade da tarefa; explicação das consequências de alterações; interpretação de representações; e discussão de possibilidades ou limites de decisões automatizadas. Frequências e transições entre níveis da rubrica ajudarão a descrever o conjunto, sem atribuir à plataforma um efeito causal isolado.
   
   \textbf{Palavras-chave}: Ensino de Computação; Machine Learning; Ensino Médio; Sequência Didática; Plataforma Interativa; Produto Educacional.

 \end{otherlanguage*}
\end{resumo}

% O resumo em inglês (Abstract) será incluído após a validação da versão portuguesa.
